% Standalone problem document, generated from the adjacent README.md.
% Compile with XeLaTeX. Mathematical content is maintained in Markdown.
\documentclass[11pt,a4paper]{article}
\usepackage[fontsize=10.5pt]{scrextend}
\usepackage[margin=22mm,headheight=15pt,headsep=9mm,footskip=11mm]{geometry}
\usepackage{fontspec}
\setmainfont{texgyrepagella}[Extension=.otf,UprightFont=*-regular,BoldFont=*-bold,ItalicFont=*-italic,BoldItalicFont=*-bolditalic]
\setsansfont{texgyreheros}[Extension=.otf,UprightFont=*-regular,BoldFont=*-bold,ItalicFont=*-italic,BoldItalicFont=*-bolditalic]
\setmonofont{texgyrecursor}[Extension=.otf,UprightFont=*-regular,BoldFont=*-bold,ItalicFont=*-italic,BoldItalicFont=*-bolditalic,Scale=MatchLowercase]
\usepackage{amsmath,amssymb}
\usepackage{unicode-math}
\setmathfont{texgyrepagella-math.otf}
\usepackage{xcolor,microtype,enumitem,needspace,fancyhdr}
\usepackage{bookmark}
\definecolor{ink}{HTML}{17344B}
\definecolor{accent}{HTML}{197D80}
\definecolor{muted}{HTML}{596773}
\hypersetup{colorlinks=true,linkcolor=accent,urlcolor=accent,pdftitle={IE-19 - Sharp
inverse norm bound from upper bounds on matrix
entries},pdfauthor={Open Problems in Numerical Linear Algebra}}
\urlstyle{same}
\setlength{\parindent}{0pt}
\setlength{\parskip}{5pt plus 1pt minus 1pt}
\linespread{1.04}
\setlength{\emergencystretch}{3em}
\widowpenalty=10000
\clubpenalty=10000
\displaywidowpenalty=10000
\allowdisplaybreaks[1]
\setlist{nosep,leftmargin=1.5em}
\providecommand{\tightlist}{\setlength{\itemsep}{0pt}\setlength{\parskip}{0pt}}
\providecommand{\pandocbounded}[1]{#1}
\pagestyle{fancy}
\fancyhf{}
\fancyhead[L]{\sffamily\footnotesize\color{muted}OPEN PROBLEMS IN NUMERICAL LINEAR ALGEBRA}
\fancyhead[R]{\sffamily\bfseries\footnotesize\color{accent}IE-19}
\fancyfoot[L]{\parbox[b]{0.9\textwidth}{\sffamily\fontsize{8}{10}\selectfont\color{muted}Editorial ratings; a bounded search cannot certify that a problem remains unsolved.\\Literature check: 2026-09-12}}
\fancyfoot[R]{\sffamily\footnotesize\color{muted}\thepage}
\renewcommand{\headrulewidth}{0.3pt}
\renewcommand{\headrule}{\hbox to\headwidth{\color{accent}\leaders\hrule height \headrulewidth\hfill}}
\makeatletter
\renewcommand{\section}{\Needspace{5\baselineskip}\@startsection{section}{1}{0pt}{12pt plus 2pt minus 2pt}{4pt}{\sffamily\bfseries\large\color{ink}}}
\renewcommand{\subsection}{\Needspace{4\baselineskip}\@startsection{subsection}{2}{0pt}{9pt plus 1pt minus 1pt}{3pt}{\sffamily\bfseries\normalsize\color{ink}}}
\makeatother
\setcounter{secnumdepth}{0}
\begin{document}
{\sffamily\small\color{muted}Linear systems and elimination\par}
\vspace{3pt}
{\raggedright\hyphenpenalty=10000\exhyphenpenalty=10000\sffamily\bfseries\fontsize{21}{24}\selectfont\color{ink}Sharp
inverse norm bound from upper bounds on matrix entries\par}
\vspace{7pt}
{\sffamily\small\textbf{Difficulty:} challenging\quad\textbf{Importance:} interesting
to specialist\par}
{\sffamily\small\bfseries\color{accent}Status: Lean verified\par}
\vspace{4pt}
{\color{accent}\hrule height 0.8pt}
\vspace{6pt}
\textbf{Topic:} Conditioning of positive diagonally dominant linear
systems.

\subsection{Lean proof and verification evidence ---
2026-09-12}\label{lean-proof-and-verification-evidence-2026-09-12}

\textbf{The original IE-19 conjecture is false, with a complete
Lean-verified counterexample.} The
\href{https://github.com/sgstepaniants/OpenProblemsInNLA/tree/531941ca0062049ccf03d3f2df418ea805ac4036/linear-systems-and-elimination/IE-19/lean}{proof
at revision 531941c} proves that the admissible matrix with diagonal
entries \(2\) and off-diagonal entries \(1/2\), at \(n=3\) and
\(m=\alpha=1\), has a genuine inverse with infinity norm \(7/9<5/4\).
The formal statement keeps the original entrywise bounds, weak diagonal
dominance and all parameter quantifiers, without adding an invertibility
premise. This counterexample refutes the full original lower-bound and
sharp conjectures. The manuscript's additional sharp-infimum and
nonattainment theorem remains informally reviewed and is outside this
Lean certificate.

\textbf{Mathematical proof:} Matthew J. Colbrook, Department of Applied
Mathematics and Theoretical Physics, University of Cambridge.
\textbf{Lean formalization:} George Stepaniants, Department of Computing
and Mathematical Sciences, California Institute of Technology, Pasadena,
California, USA, with AI-agent assistance.

The checked declarations in
\href{https://github.com/sgstepaniants/OpenProblemsInNLA/blob/531941ca0062049ccf03d3f2df418ea805ac4036/linear-systems-and-elimination/IE-19/lean/Solution.lean}{Solution.lean}
are:

\begin{itemize}
\tightlist
\item
  \textit{NLA.IE19.counterexample}
\item
  \textit{NLA.IE19.not\_lowerBoundConjecture}
\item
  \textit{NLA.IE19.not\_sharpConjecture}
\end{itemize}

Two independent agents
\href{https://github.com/ajt60gaibb/OpenProblemsInNLA/blob/main/linear-systems-and-elimination/IE-19/lean/reviews}{reviewed
the complete proof} against the original target.
\href{https://github.com/sgstepaniants/OpenProblemsInNLA/actions/runs/34703188616}{Run
34703188616} executed Comparator remotely on GitHub Actions Ubuntu
24.04, matched all three formal statements and replayed the solution
through Lean's default kernel. The
\href{https://github.com/ajt60gaibb/OpenProblemsInNLA/blob/main/linear-systems-and-elimination/IE-19/lean/verification/linux-2026-09-12}{archived
log and operational review} record the actual run, exact source hashes,
original artifact digests, and executed rejection controls for
\textit{sorry} and native-execution axioms. The
\href{https://github.com/ajt60gaibb/OpenProblemsInNLA/blob/main/linear-systems-and-elimination/IE-19/lean/verification/axioms.log}{transitive
axiom report} contains only \textit{propext}, \textit{Classical.choice},
and \textit{Quot.sound}. The independent operational referee checked the
downloaded evidence locally; it did not run Linux on the local macOS
machine. The reviews are independent agent reviews, not external human
peer review.

The proof pins \textbf{Lean 4.33.1},
\href{https://github.com/alerad/leancert/tree/621a43d7cf21f87872392a01e874f2f1dbddc926}{LeanCert
621a43d} and
\href{https://github.com/leanprover-community/mathlib4/tree/0df444a360eaa60ab8c11dca51a86af692955474}{Mathlib
0df444a}. The
\href{https://github.com/ajt60gaibb/OpenProblemsInNLA/blob/main/linear-systems-and-elimination/IE-19/lean/formalization.yaml}{formalization
manifest} and
\href{https://github.com/ajt60gaibb/OpenProblemsInNLA/blob/main/linear-systems-and-elimination/IE-19/lean/NUMERICAL_TARGETS.md}{numerical
targets} describe scope and dependencies. From a checkout of the
verified revision, with the documented
\href{https://github.com/ajt60gaibb/OpenProblemsInNLA/blob/main/tools/lean/HARNESS.md}{non-root
Linux prerequisites}, reproduce the check with:

\begin{verbatim}
tools/lean/bootstrap.sh /absolute/path/to/nla-lean-tools
tools/lean/selftest.sh /absolute/path/to/nla-lean-tools
tools/lean/verify.sh \
  linear-systems-and-elimination/IE-19/lean \
  /absolute/path/to/nla-lean-tools
\end{verbatim}

\subsection{Independently reviewed resolution -
2026-09-11}\label{independently-reviewed-resolution---2026-09-11}

\textbf{Negative resolution and sharp replacement.} Section 1 gives an
admissible positive symmetric strictly diagonally dominant \(3\times3\)
matrix with inverse infinity norm \(7/9\), below the proposed comparison
value \(5/4\). Theorem 1 in Section 2 proves that the exact infimum over
the displayed class is \(1/(\alpha+m)\) for every allowed parameter
choice, and strict positivity prevents attainment. The order is
entrywise; stronger comparisons of dominance margins are outside the
result.

\textbf{Author:} Matthew J. Colbrook, Department of Applied Mathematics
and Theoretical Physics, University of Cambridge.
\href{https://github.com/ajt60gaibb/OpenProblemsInNLA/blob/main/references/colbrook-recovered-2026-09-11/manuscripts/IE-19.pdf}{Complete
manuscript},
\href{https://github.com/ajt60gaibb/OpenProblemsInNLA/blob/main/references/colbrook-recovered-2026-09-11/verification/reviews/IE-19-review.md}{independent
proof review}, and
\href{https://github.com/ajt60gaibb/OpenProblemsInNLA/blob/main/references/colbrook-recovered-2026-09-11/README.md}{submission
record}. The supplied notes were reconstructed with substantial AI
assistance. The 2026-09-11 review was independent agent verification,
without external human peer review or formal proof-assistant
certification. The later Lean verification above covers the complete
negative resolution of the original conjecture. No novelty or priority
claim is made.

The difficulty, importance and rating rationale below are historical
assessments of the original open target. Original statements, references
and dated audits are preserved.

\textbf{Rating rationale:} Challenging reflects a sharp inverse-norm
inequality over a constrained matrix family; specialist impact concerns
extremal conditioning for positive diagonally dominant systems.

\subsection{Problem statement}\label{problem-statement}

Let \(n\ge3\), \(m>0\), and \(\alpha\ge(n-2)m\), and define

\[
S=\alpha I_n+m\mathbf1\mathbf1^T,
\qquad \mathbf1=(1,\ldots,1)^T\in\mathbb R^n.
\]

For every real symmetric matrix \(J\) satisfying

\[
0<J_{ij}\le S_{ij}\quad(1\le i,j\le n),
\qquad J_{ii}\ge\sum_{j\ne i}J_{ij}\quad(1\le i\le n),
\]

does the sharp bound

\[
\|J^{-1}\|_\infty\ge
\frac{\alpha+2m(n-1)}{\alpha(\alpha+mn)}
=\|S^{-1}\|_\infty
\]

hold, with equality if and only if \(J=S\)? Here
\(\|M\|_\infty=\max_i\sum_j|M_{ij}|\). The hypotheses make \(J\)
nonsingular. The inequalities on entries are entrywise, not Loewner
inequalities.

\subsection{Why it matters}\label{why-it-matters}

This would give an optimal lower bound on inverse amplification for a
structured family of linear systems using readily available entry
bounds.

\newpage
\subsection{References}\label{references}

\begin{itemize}
\tightlist
\item
  C. J. Hillar, S. Lin, and A. Wibisono,
  \href{https://arxiv.org/abs/1203.6812}{Inverses of symmetric,
  diagonally dominant positive matrices and applications}, §8,
  Conjecture 8.1, p.~17; the opening notation defines entrywise
  ordering. The parameter range for \(S\) is inherited from Theorem 6.1
  and the paragraph preceding Conjecture 8.1.
\item
  C. J. Hillar and A. Wibisono,
  \href{https://redwood.berkeley.edu/wp-content/uploads/2018/01/hillar2015hadamard.pdf}{A
  Hadamard-type lower bound for symmetric diagonally dominant positive
  matrices}, \emph{Linear Algebra and its Applications} 472 (2015),
  135--141, §3, Conjecture 3.4 and Theorem 3.5. This follow-up settles
  the separate determinant conjecture, not the inverse norm claim.
\end{itemize}

\subsection{Earlier status check ---
2026-09-08}\label{earlier-status-check-2026-09-08}

On 2026-09-08, searched the original title and combinations of
``Hillar'', ``Lin'', ``Wibisono'', ``Conjecture 8.1'', ``Conjecture
7.1'', ``Tight bounds on the infinity norm'', and ``proof''. The same
inverse claim also occurs as Conjecture 7.1 in the authors' revised
manuscript titled \emph{Tight bounds on the infinity norm of inverses of
symmetric diagonally dominant positive matrices}. No solution was
located. The 2015 determinant resolution was checked separately. No
recent explicit reaffirmation was found. The displayed parameter
restriction makes explicit the diagonally dominant comparison matrix
used by the source; it does not replace the conjecture's entrywise
comparison by the stronger bound on diagonal dominance margins mentioned
before it.

\subsection{Audit update --- 2026-09-10}\label{audit-update-2026-09-10}

Rechecked \href{https://arxiv.org/pdf/1203.6812}{Hillar--Lin--Wibisono},
§8, Conjecture 8.1: the extremizer and entrywise hypothesis agree with
this statement. Searches for that conjecture and later inverse bounds
found no proof of the displayed sharp inequality. Later determinant
results cited above concern a different objective. The open verdict
relies on this historical explicit conjecture and the bounded follow-up
search.
\end{document}
